%% =====================================================================
%%  The formalization companion: Lean representations and scope.
%%
%%  GENERATED FILE -- do not edit.  Produced by make_standalone.py from
%%  csp-blueprint.tex and the per-section sources; edit those and regenerate.
%%  Line references in a review should therefore be resolved back to the
%%  section sources, whose names appear in the banners below.
%%  Generated from revision 3434607 on 2026-08-02.
%% =====================================================================

%% ---------------------------------------------------------------- csp-blueprint.tex
\documentclass[11pt]{article}
\def\hazardheading{Formalization note}
\def\docheading{CSP dichotomy --- blueprint}
%% ---------------------------------------------------------------- preamble.tex
% Shared preamble for the three CSP dichotomy documents: the proof paper
% (csp-proof.tex), the source audit (csp-audit.tex), and the formalization
% blueprint (csp-blueprint.tex).
%
% House style follows the earlier `zhuk_centers` blueprint: numbered theorem
% environments shared across a single counter, a `longtheorem` upright variant
% for statements with many numbered hypotheses, and explicit `convention`
% environments for standing hypotheses.
%
% Each main file may set \hazardheading and \docheading before \input-ing this
% file; the defaults below apply otherwise.
\usepackage[a4paper,margin=27mm,headheight=14pt]{geometry}
\usepackage[T1]{fontenc}
\usepackage{lmodern}
\usepackage{amsmath,amssymb,amsthm}
\usepackage{longtable,array}
\usepackage{fancyhdr}
\usepackage[hidelinks]{hyperref}

\newtheoremstyle{mainplain}{9pt}{9pt}{\itshape}{}{\bfseries}{.}{0.55em}{}
\newtheoremstyle{mainroman}{9pt}{9pt}{\normalfont}{}{\bfseries}{.}{0.55em}{}
\theoremstyle{mainplain}
\newtheorem{theorem}{Theorem}[section]
\newtheorem{lemma}[theorem]{Lemma}
\newtheorem{proposition}[theorem]{Proposition}
\newtheorem{corollary}[theorem]{Corollary}
\theoremstyle{mainroman}
\newtheorem{longtheorem}[theorem]{Theorem}
\newtheorem{longlemma}[theorem]{Lemma}
\newtheorem{longproposition}[theorem]{Proposition}
\newtheorem{longcorollary}[theorem]{Corollary}
\theoremstyle{definition}
\newtheorem{definition}[theorem]{Definition}
\newtheorem{remark}[theorem]{Remark}
\newtheorem{convention}[theorem]{Convention}
\providecommand{\hazardheading}{Formalization note}
\newtheorem{hazard}[theorem]{\hazardheading}
\newtheorem{entry}[theorem]{Entry}

\setlength{\parindent}{1.25em}\setlength{\parskip}{0pt}\setlength{\emergencystretch}{2em}
\newcommand{\alphlabels}{\renewcommand{\labelenumi}{\textup{(\alph{enumi})}}%
  \renewcommand{\theenumi}{\alph{enumi}}}
\newcommand{\romanlabels}{\renewcommand{\labelenumi}{\textup{(\roman{enumi})}}%
  \renewcommand{\theenumi}{\roman{enumi}}}
\providecommand{\tightlist}{\setlength{\itemsep}{0pt}\setlength{\parskip}{0pt}}
\numberwithin{equation}{section}
\setcounter{tocdepth}{2}\setcounter{secnumdepth}{3}
\pagestyle{fancy}\fancyhf{}
\providecommand{\docheading}{CSP dichotomy}
\fancyhead[L]{\small\docheading}\fancyhead[R]{\small\nouppercase{\leftmark}}
\fancyfoot[C]{\thepage}\renewcommand{\headrulewidth}{0.3pt}

% ------------------------------------------------------- cross-document keys
% The three documents are compiled separately, so a reference from one to a
% statement of another cannot be a \ref.  It is the statement's *label*, set in
% typewriter, which is stable and mechanically checkable.
\newcommand{\pkey}[1]{\texttt{#1}}
% A zero-width break opportunity, for the source's very long statement names.
\newcommand{\brk}{\discretionary{}{}{}}
\newcommand{\inproof}[1]{the proof paper, \pkey{#1}}

% ---------------------------------------------------------------- operators
\DeclareMathOperator{\Sg}{Sg}
\DeclareMathOperator{\Clo}{Clo}
\DeclareMathOperator{\Pol}{Pol}
\DeclareMathOperator{\Inv}{Inv}
\DeclareMathOperator{\Var}{Var}
\DeclareMathOperator{\Con}{Con}
\DeclareMathOperator{\Sol}{Sol}
\DeclareMathOperator{\Expanded}{Exp}
\DeclareMathOperator{\CSP}{CSP}
\DeclareMathOperator{\GenSAT}{SAT}
\DeclareMathOperator*{\proj}{pr}
\newcommand{\ConOne}{\operatorname{Con}_1}
\DeclareMathOperator{\LeftLinked}{LeftLinked}
\DeclareMathOperator{\RightLinked}{RightLinked}

% ---------------------------------------------------------------- algebras
\newcommand{\bA}{\mathbf A}
\newcommand{\bB}{\mathbf B}
\newcommand{\bC}{\mathbf C}
\newcommand{\bD}{\mathbf D}
\newcommand{\bE}{\mathbf E}
\newcommand{\bR}{\mathbf R}
\newcommand{\bS}{\mathbf S}
\newcommand{\bZ}{\mathbf Z}
\newcommand{\Vn}{\mathcal V_{n}}

% ---------------------------------------------------------- subuniverse types
% The six types of Zhuk's theory.  Typed as upright small capitals so that a
% type never reads as a product of variables.
\newcommand{\TBA}{\textup{\textsc{ba}}}
\newcommand{\TC}{\textup{\textsc{c}}}
\newcommand{\TS}{\textup{\textsc{s}}}
\newcommand{\TD}{\textup{\textsc{d}}}
\newcommand{\TL}{\textup{\textsc{l}}}
\newcommand{\TPC}{\textup{\textsc{pc}}}
\newcommand{\MT}[1]{\mathcal M#1}
\newcommand{\TML}{\MT\TL}
\newcommand{\TMPC}{\MT\TPC}
\newcommand{\TMD}{\MT\TD}

% `C <_{T}^{A} B' and its relatives.
\newcommand{\subt}[2][]{<_{#2}^{#1}}
\newcommand{\subte}[2][]{\le_{#2}^{#1}}
\newcommand{\dsubt}[2][]{\dot<_{#2}^{#1}}
\newcommand{\dsubte}[2][]{\dot\le_{#2}^{#1}}
\renewcommand{\lll}{\mathrel{\lhd\!\lhd\!\lhd}}
\newcommand{\dlll}{\mathrel{\dot{\lhd}\!\dot{\lhd}\!\dot{\lhd}}}
\newcommand{\sdle}{\le_{\mathrm{sd}}}

% congruence cover  sigma^*
\newcommand{\cover}[1]{#1^{*}}
% bridge collapse
\newcommand{\bcol}[1]{\widetilde{#1}}
% absorption
\newcommand{\ab}{\lhd}
\newcommand{\abin}{\lhd_{\mathrm{bin}}}
\newcommand{\abc}{\lhd_{\mathrm{c}}}

\hypersetup{pdftitle={The CSP Dichotomy Theorem: a formalization blueprint},
  pdfsubject={Lean representations, missing Mathlib layers, order of attack},
  pdfauthor={}}

\title{The CSP Dichotomy Theorem\\[2pt]
  \large A formalization blueprint}
\author{}\date{}

\begin{document}
\maketitle\vspace{-2.2em}
\begin{abstract}
This is the formalization companion to \emph{The CSP Dichotomy Theorem}, a
corrected exposition of Zhuk's simplified proof \cite{ZhukSimplified}. It
carries the material that is about a Lean~4 development rather than about the
mathematics: which of the five statements a formalization would be proving and
what each costs (\S\ref{sec:bp-targets}--\S\ref{sec:complexity-layer}), the
representation decisions and their reasons (\S\ref{sec:design}), the layers
Mathlib does not supply (\S\ref{sec:mathlib-gaps}), the order in which the
theorems should be attacked (\S\ref{sec:order}), and what governs the schedule
(\S\ref{sec:scope}).

References of the form \texttt{lem:\dots} are labels of statements in the proof
paper, which is compiled separately; the two documents share no cross-reference
machinery on purpose, so that neither can be read as evidence for the other.
Claims here about the state of a repository are claims about a repository, and
should be checked against a commit hash rather than believed.
\end{abstract}

\tableofcontents\clearpage

%% ---------------------------------------------------------------- bp1-complexity.tex
\section{What a formalization can claim}\label{sec:bp-targets}

The dichotomy theorem, as stated in the proof paper, is conditional on a model
of computation: ``solvable in polynomial time'' quantifies over an encoding of
instances as strings, a machine, and a polynomial bound on its step count, and
``\textup{NP}-complete'' additionally quantifies over all problems in
\textup{NP} and over polynomial-time many-one reductions. Mathlib supplies
Turing machines (\texttt{Mathlib/Computability/TuringMachine/}) and a notion of
computability, but it has \emph{no} complexity class, no polynomial-time
predicate, and no notion of polynomial-time reduction. So a formalization must
say which of the statements it is proving, and the proof paper's five layers
are the answer.

\subsection{Where the work actually is}\label{sec:program-axis}

An earlier draft of this section concluded from that gap that the
complexity-theoretic half should be deferred because building it would dominate
the project. That inference was wrong, and it is worth recording why, because
the corrected version is what organises the rest of this document.

The complexity vocabulary itself --- languages, a cost model, $\mathrm P$,
$\mathrm{NP}$ by verifiers, $\le_{p}$, \textup{NP}-hardness,
\textup{NP}-completeness, and a generic \textup{NP}-hard problem --- has since
been built, in about $850$ lines. That is consistent with the two existing
Cook--Levin formalizations, whose class layers are $702$ lines (Isabelle) and
$992$ lines (Coq) against totals of $56\,514$ and $\approx 16\,500$. Defining
the classes is cheap in every development that has done it.

What is expensive is not complexity theory. It is any statement that requires
one to \emph{exhibit a program}. In a formalization, ``$f$ is computable in
polynomial time'' is not a property one observes of a mathematical function; it
is an existential over programs, discharged only by writing one and proving a
bound on its step count. The dividing line runs along that axis, and it cuts
across the algebra/complexity split rather than agreeing with it:

\begin{center}
\begin{tabular}{@{}lll@{}}
\hline
Target & Exhibit a program? & Cost \\
\hline
\textup{(T0)} the algebraic core & no & the mathematics \\
\textup{(T1)} $\textsc{Solve}$ is correct & no (but see Formalization note~\ref{haz:t1-computable}) & moderate \\
\textup{(T2)} $\textsc{Solve}$ is polynomial & \textbf{yes --- the largest one here} & \textbf{large} \\
\textup{(H0)} $\Gamma$ pp-interprets \textsc{Nae} & no & the mathematics \\
\textup{(H1)} pp-interpretation $\Rightarrow\le_{p}$ & yes, a syntactic substitution & small \\
\hline
\end{tabular}
\end{center}

The blueprint is written for \textup{(T0)}, \textup{(T1)} and \textup{(H0)}.
\textup{(H1)} is a small addition on top; \textup{(T2)} is neither small nor
mathematics, but it is not \emph{unrelated} to the dichotomy either --- it is
the assertion that the algorithm one has just proved correct is also fast, and
it is the second-largest component of the project.

\begin{hazard}[\textup{(T1)} is vacuous unless $\textsc{Solve}$ computes]
\label{haz:t1-computable}
``$\textsc{Solve}$ written as a total function'' hides a requirement that no
cost model would catch. If $\textsc{Solve}$ is defined using classical choice
--- for instance by deciding ``does $D_{x}$ have a nonempty proper binary
absorbing subuniverse?'' with \texttt{Classical.dec} --- then
$\textsc{Solve}(\mathcal I)=\texttt{true}\iff\mathcal I$ has a solution is a
true theorem that says nothing algorithmic, because $\textsc{Solve}$ does not
reduce. To mean what it appears to mean, \textup{(T1)} needs every predicate
the algorithm branches on to carry a genuine decision procedure.

That is a real obligation with a locatable cost. Absorption is defined by an
existential over \emph{terms} --- an infinite type --- so ``$C$ is a binary
absorbing subuniverse of $B$'' is not decidable on its face. It becomes
decidable through a bound on the arity of a sufficient witness, which is a
theorem and not a definition. For the types actually used here the bound is
cheap and should be taken that way rather than from a general absorption-arity
theorem: $\TBA$ has a binary witness by definition, and central absorption has
a ternary witness by \inproof{lem:central-ternary}. Only finitely many binary
and ternary operations then need be checked, together with finitely many
congruences and subsets. The same remark applies to the searches inside
$\textsc{SolveLinear}$.
\end{hazard}

\begin{hazard}[\textup{(T2)} is not a wrapper]\label{haz:t2-is-not-small}
Two things make \textup{(T2)} much larger than the phrase ``complexity
wrapper'' suggests.

\emph{The program is enormous.} \textup{(T2)} requires $\textsc{Solve}$ ---
\textsc{ForceConsistency}, the search for a strong subuniverse,
\textsc{SolveLinear}, and the mutual recursion between them --- to be expressed
in a cost model, with a proved polynomial bound. Every existing measurement
says this kind of work dominates: it is the $50\,000$ lines of Turing-machine
engineering in the Isabelle Cook--Levin development, against $702$ for the
classes themselves.

\emph{The statement is non-uniform.} Because $\Gamma$ is fixed, the polynomial
may depend on $\Gamma$, and the searches over subsets of $A$, over congruences
on $A$, and over term operations of bounded arity are all constant-time. But
``constant'' here means the search is compiled away into a program whose
\emph{size} depends on $|A|$. So $\textsc{Solve}$ is not one program: it is a
family $\textsc{Solve}_{\Gamma}$ with a polynomial $p_{\Gamma}$ for each
$\Gamma$, and \textup{(T2)} must be stated in that form. Attempting a bound
uniform in $\Gamma$ would be proving something else, and something false.
\end{hazard}

\begin{hazard}[From pseudocode to a Lean function]
\label{haz:algorithm-formalization}
\begin{enumerate}\alphlabels\tightlist
\item \emph{Termination is not structural.} $\textsc{Solve}$ recurses on
  strictly smaller domains, and $\textsc{SolveLinear}$ loops on a strictly
  decreasing dimension. Neither is a structural recursion; both need an
  explicit well-founded measure, and the two interleave ---
  $\textsc{SolveLinear}$ calls $\textsc{Solve}$ on a smaller domain, which
  calls $\textsc{SolveLinear}$ again. The measure is lexicographic:
  $\bigl(\sum_{x}|D_{x}|,\ \dim\phi\bigr)$.
\item \emph{Several steps quantify over objects that are only finitely many by
  accident.} ``If some $D_{x}$ has a strong subset $B$'' quantifies over
  subsets of $D_{x}$, over congruences on $D_{x}$, and --- through absorption
  --- over \emph{term operations}, of which there are infinitely many. See
  Formalization note~\ref{haz:t1-computable} for the route to decidability that does not
  depend on an uncited general arity bound.
\item \emph{Polynomial time is a different subject.} The bound in
  \cite{ZhukJACM} counts calls and bounds the recursion depth by $|A|$; making
  that formal needs a cost model. See \S\ref{sec:complexity-layer}.
\end{enumerate}
\end{hazard}

\section{The complexity layer}\label{sec:complexity-layer}

Targets \textup{(T2)} and \textup{(H1)} are stated here and not proved. The
point of stating them is to be precise about what is \emph{not} being claimed.

\begin{definition}[What a cost model would have to supply]\label{def:cost-model}
To state \textup{(T2)} one needs: an encoding of instances of $\CSP(\Gamma)$ as
strings; a machine model with a step count; a predicate ``$f$ is computed in
time bounded by a polynomial in the input length''; and closure of that
predicate under composition. To state \textup{(H1)} one needs, in addition, a
polynomial-time many-one reduction relation $\le_{p}$ and the class
\textup{NP}.
\end{definition}

Mathlib supplies the machine (\texttt{Turing.TM2}), an encoding
(\texttt{Computability.Encoding}), and a time-bounded computation predicate
(\texttt{Turing.TM2ComputableInPolyTime}). It does not supply a complexity
\emph{class}, \textup{P}, \textup{NP}, \textup{NP}-completeness, $\le_{p}$, or
\textsc{Sat}; and the composition lemma
\texttt{TM2ComputableInPolyTime.comp} is an open \texttt{proof\_wanted} in
Mathlib itself. Moreover \texttt{TM2ComputableInPolyTime} applies to total
functions $\alpha\to\beta$, never to languages, so even the existing predicate
would need reshaping.

\begin{remark}[This layer now exists, and it was cheap]\label{rem:complexity-scope}
An earlier draft of this remark said that building the complexity layer was
``a substantial project --- comparable in size to the algebraic core'' and that
attempting it as part of this project ``would be a scoping error''. That was
wrong, and it has been tested rather than argued: the vocabulary --- words, a
cost model, $\mathrm P$, $\mathrm{NP}$ by verifiers, $\le_{p}$,
\textup{NP}-hardness, \textup{NP}-completeness, and a generic
\textup{NP}-hard problem --- is about $850$ lines of Lean, \texttt{sorry}-free.

So the reason to keep \textup{(T2)} out of scope is \emph{not} that complexity
theory is unavailable. It is Formalization note~\ref{haz:t2-is-not-small}: \textup{(T2)}
requires exhibiting $\textsc{Solve}$ as a program with a proved step bound, and
that is the largest program-construction task in the project. \textup{(H1)},
by contrast, needs only a syntactic substitution and should be built.
\end{remark}

\begin{remark}[What the exercise cost, and what it caught]
\label{rem:complexity-defects}
Two defects surfaced in those $850$ lines, both in \emph{statements} rather
than proofs, and neither of a kind a type-checker reports. The generic problem
was first defined with a single budget bounding both the certificate length and
the running time, which makes the theorem false --- the two directions pull the
bound opposite ways. And the language was first given only the tests ``register
is empty'' and ``register's first bit is $1$'', with which no loop that
consumes its input can be written at all; that was found by trying to write the
first program.

The ratio is worth noting for the rest of this project. Two statement-level
errors in $850$ lines, against fourteen found in Zhuk's fifty pages by readers
who did not formalize. The defect list of the audit document should be expected
to grow under formalization, not to be complete.
\end{remark}

\begin{remark}[\textup{NP} membership]\label{rem:np-membership}
The proof paper's hard half concludes \textup{NP}-hardness. Recovering
\textup{NP}-completeness needs the standard theorem that every fixed
finite-template CSP lies in \textup{NP} under the chosen encoding: a solution
is a certificate of size linear in the number of variables, and checking it
means evaluating each constraint, which for fixed $\Gamma$ is a table lookup.
It is small, it needs the cost model, and it belongs with \textup{(H1)}.
\end{remark}
\clearpage
%% ---------------------------------------------------------------- bp2-lean.tex
\section{The Lean development}\label{sec:formalization}

\subsection{What is reused}\label{sec:reuse}

The predecessor development \texttt{ZhukLean} is a complete, \texttt{sorry}-free
Lean~4 formalization of Zhuk's centre theorem, built on
\texttt{Mathlib.ModelTheory}. It is consumed here as a \texttt{lake} dependency,
not copied. It supplies:

\begin{itemize}\tightlist
\item products of structures (\texttt{piStructure}, \texttt{prodStructure},
  coordinatewise realization, reindexing and evaluation homomorphisms) --- the
  one part of the universal-algebra vocabulary Mathlib lacks;
\item absorption in the tuple form (\texttt{Witnesses}, \texttt{Absorbs},
  \texttt{BinAbsorbs}), Taylor terms (\texttt{IsTaylorOn}), and central
  absorption (\texttt{CentrallyAbsorbs});
\item the relational description of absorption by essential relations
  (Barto--Kazda), which is what converts absorption at some arity into
  absorption at arity~$3$;
\item Zhuk's centre theorem itself (\texttt{zhuk\_center}) and the ternary
  collapse (\texttt{exists\_ternary\_\brk witnesses}) --- which is precisely
  \inproof{lem:central-ternary}, cited without proof in \cite{ZhukSimplified}.
\end{itemize}

Those also discharge, in advance, every ingredient of \cite[Lem.~6.11]{ZhukStrong},
which the proof paper's stable-intersection theorem needs in its type-$\TC$
case: essentiality, the Barto--Kazda criterion, and ``central implies ternary
absorbing''.

\subsection{What is built here}\label{sec:built}

\begin{center}
\begin{tabular}{@{}llr@{}}
\hline
Module & Contents & Status \\
\hline
\texttt{CSP/Product} & products of structures & complete \\
\texttt{CSP/Congruence} & \texttt{Congruence L M}, order, \texttt{sInf}, quotient & complete \\
\texttt{CSP/Irreducible} & irreducibility, existence and uniqueness of $\cover\sigma$ & complete \\
\texttt{CSP/Bridge} & bridges, collapse, composition, linear/PC & 2 imports open \\
\texttt{CSP/Types} & the six types, multi-types, $\lll$ & complete \\
\texttt{CSP/Instance} & instances, reductions, consistency, cruciality & complete \\
\texttt{CSP/Main} & the main statements & targets \\
\hline
\end{tabular}
\end{center}

Everything below \texttt{CSP/Main} is \texttt{sorry}-free. The two open items in
\texttt{CSP/Bridge} are \inproof{lem:bridge-comp} and its collapse identity,
both imported from \cite{ZhukJACM}. These are claims about a repository, not
about the mathematics; they should be read against a commit hash, and they are
deliberately kept out of the proof paper for that reason.

\subsection{Design decisions}\label{sec:design}

\begin{description}
\item[Build on \texttt{Mathlib.ModelTheory}, keep the language generic.]
  A bespoke ``finite algebra with one WNU operation'' type would bundle the
  standing hypotheses as fields and avoid threading them, which is a real cost
  in the predecessor development. But it would discard
  \texttt{Substructures.lean} --- 985 lines of exactly the $\Sg$ API needed,
  including \texttt{mem\_closure\_iff\_exists\_term} with the variable type
  equal to the generating set, an interface no hand-rolled clone matches --- and
  it would discard the predecessor's 1600 lines wholesale. Congruences and
  quotients are also \emph{cheaper} in \texttt{ModelTheory}, not dearer.
  Instantiating to $\Vn$ is a thin layer on top. Carry the standing hypotheses
  as typeclass instances on the structure rather than as ambient hypotheses, so
  that the audit of which result consumes which hypothesis is mechanical.

\item[Terms and term operations are different types.]
  \cite{ZhukSimplified} writes $t\in\Clo(\mathbf A)$ and treats $t$ as both a
  syntactic term and the operation it induces. Every statement of the proof
  paper that quantifies over $\Clo(\mathbf A)$ is about \emph{operations}, so
  the syntactic layer may be quotiented away --- but it cannot be skipped,
  because $\Sg$ is characterised through terms and because the arity
  bookkeeping in the absorption arguments is syntactic.

\item[$\lll$ is a witnessed chain, not \texttt{Relation.ReflTransGen}.]
  An earlier draft of this section said the opposite, and it was wrong: the
  proof paper's ``the congruences coming from $C\lll^{A}B$'' is a function of
  the chain and not of the pair, and two arguments turn on \emph{which} chain
  is chosen. So the primary notion is an inductive family carrying, at each
  step, the intermediate set, its type, and its witnessing congruence, with
  extension by one step as an operation on that data;
  \texttt{Relation.ReflTransGen} is its propositional truncation and may be
  used only where no congruence is named. Both are needed: the truncation
  supplies the induction principle most consumers want.

\item[Dotted relations are defined as disjunctions.]
  \texttt{DotLll C B := C = }$\varnothing$\texttt{ }$\vee$\texttt{ Lll C B}, so
  the case split is forced at every use site rather than resolved by the
  reader. Conflating $\lll$ with $\dlll$ is the commonest way to misread the
  type theory of the proof paper's \S3.

\item[Absorption is stated over coordinatewise-constrained tuples.]
  That is: for every $\bar a\in A^{k}$ with $a_{j}\in B$ for all $j\neq i$, one
  has $t(\bar a)\in B$. \emph{Not} as ``for all
  $b_{1},\dots,b_{k}\in B$ and $a\in A$, $t(b_{1},\dots,a,\dots,b_{k})\in B$''.
  The two agree, but the second phrasing
  invites the variant in which $b_{i}$ is quantified and then discarded, which
  is vacuously true when $B=\varnothing$ and at $k=1$ where the intended
  content is not vacuous. This was an actual defect in the predecessor
  blueprint.

\item[The bridge criterion is the definition of linearity.]
  Zhuk defines a linear congruence by a per-block isomorphism to
  $\mathbb Z_{p}^{m}$ and derives the bridge criterion. We reverse this: the
  criterion is first-order over finite data and is what every consumer uses;
  the block description becomes a structure theorem. See
  \inproof{haz:linear-def}.

\item[$\cover\sigma$ is data attached to irreducibility.]
  \texttt{Irreducible.exists\_isCover} produces it; there is no standalone
  function, because there is no cover without irreducibility. A pleasant
  by-product: with the definition taken literally, irreducibility \emph{implies}
  $\sigma\neq A^{2}$, since the empty family intersects to the universe.

\item[$\ConOne$ gets a neutral name until it is a congruence.]
  The raw relation is always reflexive and symmetric but not transitive; the
  notation invites the reader to assume congruence-hood. Reserve $\ConOne$ for
  the case where subdirectness and rectangularity make it one. Note also that
  it depends on the \emph{position} in the scope, not on the variable, so it is
  ill-defined if a variable occurs twice.

\item[Instances are lists; scopes are tuples; the variable type is a parameter.]
  Proofs delete a constraint and reason about $\mathcal I\setminus\{C\}$, and
  the same relation may occur twice on the same tuple inside an expanded
  covering, so \texttt{List} with membership by index is the safe choice.
  Repeated variables in a scope are allowed and occur.

\item[Coverings are built by coproduct of variable types, not by renaming.]
  \cite{ZhukSimplified} writes ``we rename the variables so that the only
  common variable of $\Upsilon_{x}$ and $\mathcal J$ is $x$''. With a concrete
  variable type this means a fresh-name supply, a renaming operation, and a
  lemma that renaming preserves each of $1$-consistency, cycle-consistency,
  cruciality, and being a covering --- eight lemmas the paper does not state.
  Making the variable type a parameter and gluing over $V_{1}\oplus V_{2}$
  quotiented by the shared variable turns all eight into transport along an
  equivalence. This is the single most important representation decision below
  the algebra, and it should be made before anything about coverings is
  written.

\item[Nothing is stubbed with \texttt{sorry} at the level of a definition.]
  \texttt{IsConnectedInstance} and \texttt{IsExpandedCovering} are
  \emph{absent} rather than defined as \texttt{sorry}, because a
  \texttt{def~:~Prop~:=~sorry} is a junk constant that downstream proofs can
  silently lean on. Absence forces the design decision; a stub hides it.
\end{description}

\subsection{What Mathlib is missing}\label{sec:mathlib-gaps}

\begin{center}\small
\begin{tabular}{@{}p{0.32\textwidth}p{0.60\textwidth}@{}}
\hline
\textbf{Layer} & \textbf{Status} \\
\hline
Products of first-order structures &
  Absent from Mathlib; supplied by \texttt{ZhukLean} (\S\ref{sec:reuse}). \\
Congruences of a structure, quotients &
  \texttt{ModelTheory} has substructures but not congruences; built here in
  \texttt{CSP/Congruence}. \\
Absorption, central subuniverses &
  Absent; supplied by \texttt{ZhukLean}. \\
Mixed-prime ``dimension'' &
  A subgroup of $\prod_{i}\mathbb Z_{q_{i}}$ with distinct $q_{i}$ is not a
  vector space. \texttt{ZMod} and \texttt{Module (ZMod p)} exist; the primary
  decomposition into $\mathbb F_{q}$-vector spaces and the codimension-one
  characterisation must be built. Needed by the proof paper's codimension-one
  theorem. \\
The Pol--Inv Galois connection &
  Absent. Roughly four printed pages, of independent value, and the single
  largest item in the \textup{(H0)} budget. \\
Complexity classes, $\le_{p}$, \textup{NP} &
  Absent; see \S\ref{sec:complexity-layer}. Built separately in $\approx 850$
  lines. \\
Polynomial completeness, Istinger--Kaiser &
  Absent. Needed only for \inproof{lem:pc-on-top}, which nothing else consumes;
  a staged development should treat ``PC congruence'' as ``irreducible and not
  linear'' throughout and prove that lemma last or not at all. \\
\hline
\end{tabular}
\end{center}

\subsection{Order of attack}\label{sec:order}

\begin{enumerate}\tightlist
\item \inproof{lem:ba-center-implies}. Small, self-contained, used everywhere,
  and adjacent to what is already formalized. Its statement needs both
  hypotheses the source drops --- subdirectness, and a common binary term ---
  and the common term should be carried in the type, as \cite{ZhukStrong}
  writes $\le_{\TBA(t)}$.
\item The mixed-prime linear algebra of \S\ref{sec:mathlib-gaps}. Independent
  of everything else; can be done in parallel.
\item \inproof{lem:bridge-comp}. Twelve source lines of relational algebra;
  closes the two open items in \texttt{CSP/Bridge}.
\item The coproduct design for coverings, then the basic properties of expanded
  coverings.
\item \inproof{thm:stable-intersection}. The one hard statement in the
  interface, and the sole source of bridges.
\item \inproof{lem:bridge-from-instance}. Eight hypotheses, ten pages.
\item \inproof{thm:main-inductive}. Last.
\end{enumerate}

\subsection{Scope, and what governs the schedule}\label{sec:scope}

The calibration point is that the predecessor formalized about $3.5$ printed
pages in about $1670$ lines of Lean --- roughly $480$ lines per page. Excluding
\S\ref{sec:complexity-layer} and \cite{ZhukSimplified}'s \S4, which is an
independent second result and not needed for the dichotomy, the transitive
closure of what must be proved is on the order of $100$--$130$ printed pages,
so $35\,000$--$60\,000$ lines of Lean.

Line count is not the binding constraint, and it should not be converted into
person-time here: any such figure would be a claim about a development process
rather than about this mathematics, and it is not evidence a reader of this
document can check.

What binds is the \emph{critical path}: the depth of the serial chain, and the
number of places where the source is wrong and new mathematics is needed. The
serial chain is
\begin{center}\small
\pkey{thm:stable-intersection}\\[2pt]
$\longrightarrow$\ \pkey{lem:bridge-from-instance}\\[2pt]
$\longrightarrow$\ \pkey{thm:main-inductive}
\end{center}
each a single large proof that cannot be split across workers --- stable
intersection is the sole source of bridges, bridge-from-instance is the only
consumer that turns them into instance-level structure, and the main induction
consumes that.

The new mathematics is a much shorter list than an earlier draft of this
section claimed, and the items on it are not the ones it named. Of the audit
document's register, thirteen defects are repaired with no new mathematics.
What remains open, and gates the formalization of \S5 of
\cite{ZhukSimplified}, is:

\begin{itemize}\tightlist
\item \inproof{lem:maximal-mult} --- the maximal multi-type extension. Its
  Case~1 is unproved and the published repair is invalid; see the audit
  document. Nothing downstream of it should be attempted first.
\item The reflexivisation of a bridge, used twice without a lemma behind it.
\item The expansion of \inproof{lem:intersection-good} from case structure to a
  full proof, and of \inproof{lem:multiply-all-linear}, which is not yet a
  proof.
\end{itemize}

Those are not throughput-limited. Schedule around them, not around the line
count.

\clearpage
%% ---------------------------------------------------------------- bibliography.tex
% Shared bibliography for the three documents.  Kept in one file so that a
% citation key means the same thing in the proof paper, the audit and the
% blueprint.
\begin{thebibliography}{99}

\bibitem{BJK} A.~Bulatov, P.~Jeavons, A.~Krokhin,
  \emph{Classifying the complexity of constraints using finite algebras},
  SIAM J. Comput.\ \textbf{34} (2005), 720--742.

\bibitem{BradyNotes} Z.~Brady, \emph{Notes on CSPs and polymorphisms},
  arXiv:2210.07383. Cited by section title and numbered statement; the
  sections used here are ``Cores and Idempotent Reducts'' and \S3.11,
  \S3.15.

\bibitem{BartoKozik} L.~Barto, M.~Kozik, \emph{Absorbing subalgebras, cyclic
  terms, and the constraint satisfaction problem}, Log. Methods Comput. Sci.
  \textbf{8} (2012), 1:07.

\bibitem{DecidingAbsorption} L.~Barto, A.~Kazda, \emph{Deciding absorption},
  Internat. J. Algebra Comput. \textbf{26} (2016), 1033--1060.

\bibitem{Bergman} C.~Bergman, \emph{Universal algebra: fundamentals and
  selected topics}, CRC Press, 2011.

\bibitem{Bulatov} A.~Bulatov, \emph{A dichotomy theorem for nonuniform CSPs},
  arXiv:1703.03021; FOCS 2017.

\bibitem{FederVardi} T.~Feder, M.~Y.~Vardi, \emph{The computational structure of
  monotone monadic SNP and constraint satisfaction}, SIAM J. Comput.\
  \textbf{28} (1999), 57--104.

\bibitem{HobbyMcKenzie} D.~Hobby, R.~McKenzie, \emph{The structure of finite
  algebras}, Contemp. Math.~76, Amer. Math. Soc., 1988.

\bibitem{IstingerKaiser} M.~Istinger, H.~K.~Kaiser, \emph{A characterization of
  polynomially complete algebras}, J. Algebra \textbf{56} (1979), 103--110.

\bibitem{MarotiMcKenzie} M.~Mar\'oti, R.~McKenzie, \emph{Existence theorems for
  weakly symmetric operations}, Algebra Universalis \textbf{59} (2008),
  463--489.

\bibitem{MinimalTaylor} L.~Barto, Z.~Brady, A.~Bulatov, M.~Kozik, D.~Zhuk,
  \emph{Minimal Taylor algebras as a common framework for the three algebraic
  approaches to the CSP}, LICS 2021.

\bibitem{RossSlides} R.~Willard, \emph{Similarity, critical relations, and
  Zhuk's bridges}, slides, AAA98 --- 98th Workshop on General Algebra,
  Dresden, 2019.
  \url{http://www.math.uwaterloo.ca/~rdwillar/documents/Slides/willard-aaa98-handout.pdf}

\bibitem{ZhukJACM} D.~Zhuk, \emph{A proof of the CSP dichotomy conjecture},
  J. ACM \textbf{67} (2020), no.~5, 1--78; arXiv:1704.01914.

\bibitem{ZhukSimplified} D.~Zhuk, \emph{A simplified proof of the CSP dichotomy
  conjecture and XY-symmetric operations}, arXiv:2404.01080v2.

\bibitem{ZhukStrong} D.~Zhuk, \emph{Strong subalgebras and the constraint
  satisfaction problem}, J. Mult.-Valued Logic Soft Comput.\ \textbf{36} (2021);
  arXiv:2005.00593.

\end{thebibliography}

\end{document}
